\documentclass[a4paper,12pt]{article}

\usepackage[utf8]{inputenc}
\usepackage{amsmath}
\usepackage{geometry}
\usepackage{listings}
\usepackage{hyperref}

\geometry{margin=2.5cm}

\title{Relazione Laboratorio HPC}
\author{}
\date{}

\begin{document}

\maketitle

% =========================================================
\section{Obiettivo dell’esperimento}

L’obiettivo del laboratorio è analizzare e confrontare le prestazioni di un sistema locale (MacBook Air con chip Apple M3) con un cluster di calcolo dell’Ateneo, valutando configurazione hardware e prestazioni teoriche in FLOPS.

% =========================================================
\section{Ambiente di calcolo}

% -----------------------------
\subsection{Ambiente locale (MacBook Air M3)}

Il sistema utilizzato è un laptop MacBook Air con chip Apple M3.

\subsubsection{Sistema operativo e architettura}

\begin{itemize}
    \item Sistema operativo: macOS 26.3.1
    \item Architettura: Apple Silicon (ARM64)
\end{itemize}

\subsubsection{CPU}

Dai test eseguiti in ambiente Docker:

\begin{itemize}
    \item Architettura rilevata: x86\_64 (emulazione container)
    \item Core: 8
    \item Socket: 1
    \item Thread per core: 1
    \item Istruzioni SIMD: ASIMD (ARM NEON), AES, SHA
\end{itemize}

L’ambiente Docker introduce una rappresentazione astratta dell’hardware fisico.

\subsubsection{GPU}

Il sistema non dispone di GPU NVIDIA, pertanto il comando \texttt{nvidia-smi} non è disponibile.

La GPU reale è:

\begin{itemize}
    \item Apple M3 GPU (10 core)
    \item Architettura: Metal 4
    \item Memoria unificata (UMA)
\end{itemize}

\subsubsection{Prestazioni di picco}

\textbf{CPU:}

\[
FLOPS_{core} = f \times 2 \times N
\]

\begin{itemize}
    \item $f = 4.06$ GHz
    \item $N = 2$ (ASIMD doppia precisione)
    \item FMA = 2
\end{itemize}

\[
FLOPS_{core} = 16.24 \text{ GFLOPS}
\]

\[
FLOPS_{CPU} = 16.24 \times 8 = 129.92 \text{ GFLOPS}
\]

\vspace{0.5cm}

\textbf{GPU:}

\begin{itemize}
    \item Frequenza stimata: 1.3 GHz
    \item Core GPU: 10
    \item FMA = 2
\end{itemize}

\[
FLOPS_{GPUcore} = 2.6 \text{ GFLOPS}
\]

\[
FLOPS_{GPU} = 26 \text{ GFLOPS}
\]

La GPU Apple utilizza API Metal e memoria unificata, senza supporto CUDA.

\subsubsection{Nota su MPI}

I nodi MPI (wn1 e wn2) sono implementati tramite container Docker sullo stesso host fisico. Essi rappresentano nodi logici e non incrementano le prestazioni hardware complessive.

% =========================================================
\subsection{Cluster di calcolo (Ateneo)}

In questa sezione verranno riportati i risultati ottenuti sul cluster di calcolo dell’Ateneo.

\subsection{Cluster di calcolo (Ateneo)}

Il nodo di calcolo dell’Ateneo è stato analizzato tramite comandi di sistema su ambiente HPC.

% =========================================================
\subsubsection{CPU (cluster)}

Dai test eseguiti, il nodo è caratterizzato da:

\begin{itemize}
    \item Architettura: x86\_64
    \item CPU: Intel Xeon Gold 6238R @ 2.20 GHz
    \item Core: 4
    \item Socket: 2 (virtualizzazione VMware)
    \item Thread per core: 1
    \item Cache L1/L2/L3: 32 KB / 1 MB / 39 MB
    \item SIMD: AVX, AVX2, AVX512
    \item FMA: presente
\end{itemize}

Il set SIMD massimo disponibile è AVX512, che corrisponde a:

\[
N = 16
\]

% -----------------------------
\paragraph{Prestazioni CPU per core}

\[
FLOPS_{core} = f \times 2 \times N
\]

dove:

\begin{itemize}
    \item $f = 2.20 \times 10^9$ Hz
    \item $N = 16$
    \item FMA = 2
\end{itemize}

\[
FLOPS_{core} = 2.20 \times 10^9 \times 2 \times 16 = 70.4 \text{ GFLOPS}
\]

% -----------------------------
\paragraph{Prestazioni CPU per nodo}

\[
FLOPS_{CPU} = 70.4 \times 4 = 281.6 \text{ GFLOPS}
\]

% =========================================================
\subsubsection{GPU (cluster)}

Durante l’esecuzione dei test la GPU non è risultata accessibile dal nodo di login.

I comandi:

\begin{verbatim}
nvidia-smi
\end{verbatim}

non sono disponibili, indicando che:
\begin{itemize}
    \item il nodo non espone GPU NVIDIA, oppure
    \item l’accesso alle GPU è riservato a nodi di calcolo dedicati
\end{itemize}

Pertanto non è stato possibile raccogliere informazioni su CUDA cores, memoria o clock GPU.

% =========================================================
\subsubsection{Prestazioni di picco (cluster)}

Sulla base della CPU analizzata:

\begin{itemize}
    \item CPU (1 core): 70.4 GFLOPS
    \item CPU (nodo): 281.6 GFLOPS
    \item GPU: non valutabile (non accessibile nel nodo corrente)
\end{itemize}
\end{document}